\chapter{Einleitung}\label{einleitung}

Am 22. Februar 2020, drei Tage nach dem rechtsextremen Anschlag von Hanau, berichtet die Tageszeitung \textsc{Die} \textsc{Welt}:

\begin{quote}
Die Terrornacht von Hanau mit elf Toten wirkt sich auch auf das Wahlverhalten der Deutschen aus. Die AfD, die nach dem rassistisch motivierten Anschlag in der Kritik steht, verlor im aktuellen RTL/ntv-Trendbarometer zwei Prozentpunkte.\footnote{Der Artikel ist unter \url{https://www.welt.de/politik/deutschland/article206062103/Neue-Forsa-Umfrage-AfD-verliert-nach-Morden-von-Hanau-spuerbar.html} {[}18.12.2025{]} online verfügbar.}
\end{quote}

Diese den Artikel einleitenden Sätze wirken aus semantischer Perspektive eher unscheinbar. Es erscheint plausibel, dass einer Partei des rechten Randes, die rassistische Positionen vertritt und in Diskurse einbringt, eine Mitschuld an einem rechtsextremen Attentat zugeschrieben wird und sie daraufhin an Stimmen einbüßt. Gleichwohl wird dieser Zusammenhang im Zitat (und im weiteren Verlauf des Artikels) sprachlich nicht explizit gemacht, sodass man ein relativ breites Wissen benötigt, um ihn zu verstehen. Nicht nur muss man wissen, dass die AfD eine rechte Partei ist -- auch dass in Demokratien Parteien gewählt werden, Menschen ihre Wahlentscheidung von aktuellen Ereignissen abhängig machen und Parteien nicht wählen, die sie als Mitverursacher von Attentaten ansehen, ist für das Verständnis notwendig. Hinter dem erwähnten „Wahlverhalten`` verbirgt sich gleich ein ganzes System von Wissenselementen: Wir müssen wissen, dass Bürger*innen eine festgelegte Anzahl an Stimmen an Parteien vergeben können, dass dies nur in einem bestimmten festgelegten mehrjährigen Rhythmus geschieht und dass die Wahltendenzen in der Zwischenzeit durch Prognosen in Form prozentualer Zustimmungswerte erfasst werden. Man muss auch wissen, dass das „RTL/ntv-Trendbarometer`` eine solche Prognose darstellt und wie viel eine Verringerung um „zwei Prozentpunkte`` dabei bedeutet, wozu es wiederum relevant ist, den üblichen Zustimmungswert der Partei zu kennen. Diese Auflistung ist keinesfalls erschöpfend, zumal jedes dieser Wissenselemente selbst wieder ein breites Feld an verstehensrelevanten Inhalten voraussetzt -- Was etwa ist eine Prognose, wie wird sie durchgeführt und welchen Stellenwert hat sie in unserer Gesellschaft?

Die Beschreibung dieses breiten Wissens und seiner sprachlichen Realisierung ist kein selbstverständlicher Teil der Linguistik. Die impliziten Anteile der Bedeutung, die sich aus dem Kontext oder dem sogenannten ‚Weltwissen` speisen, werden für gewöhnlich aus der semantischen Beschreibung ausgeklammert und entweder für die Pragmatik vorgesehen -- wo sie etwa als ‚Präsuppositionen` oder ‚Implikaturen` behandelt werden -- oder eben gar nicht beschrieben. Gerade vor Etablierung der Pragmatik als eigenständiger linguistischer (Teil\mbox{-})Disziplin fehlten diese Aspekte in der Beschreibung von Bedeutung weitenteils und man beschränkte sich stattdessen auf das Festhalten einzelner stichpunktartiger semantischer Merkmale zu Wörtern. Einer der frühesten und prominentesten Kritiker dieser reduktionistischen Ausrichtung der Semantik war der Linguist Charles J. Fillmore, der sie als „Checklist Theories of Meaning`` \autocite*[123]{fillmoreAlternativeChecklistTheories1975} bezeichnete. Seine Überlegungen zur Komplexität von Bedeutung und ihrer sprachlichen Realisierungsmöglichkeiten mündeten in die von ihm begründete Frame-Semantik. In dieser werden isolierte Beschreibungen der Semantik einzelner Wörter abgelöst durch die Bedeutungseinheit des semantischen Frames. Frames wie der von Fillmore oft als Beispiel angeführte \FN{Commercial\_event}-Frame -- oder auch das oben skizzierte System von Wissenselementen zum Wählen -- weisen Leerstellen auf, die durch einzelne lexikalische Einheiten gefüllt werden können. Eine ganze Reihe von Wörtern wie \emph{Käufer}, \emph{Verkäufer}, \emph{bezahlen}, \emph{verkaufen}, \emph{ausgeben}, \emph{Waren} oder \emph{Geld} lässt sich nur verstehen, wenn der gemeinsame, die Bedeutungsstruktur stiftende, Frame bekannt ist. Dabei wird der Frame, so Fillmore, durch die einzelnen lexikalischen Wahlen in unterschiedliche Perspektiven gesetzt: Etwas zu \emph{kaufen}, hebt die Handlung des \textsc{käufers} in Bezug auf die \textsc{ware} hervor, während etwas zu \emph{verkaufen}, den \textsc{käufer} zugunsten des \textsc{händlers} in den Hintergrund rückt. Etwas zu \emph{bezahlen} wiederum, verschiebt den Fokus auf die Rolle des \textsc{geldes} \autocite[116-117]{fillmoreFrameSemantics1982}.
\largerpage
Die Genese von Fillmores Frame-Begriff und auch dass dieser keinesfalls im luftleeren Raum entsteht, sondern u.\,a. bei der Valenzgrammatik Anleihen nimmt, arbeitet Busse \autocite*{busseFrameSemantikKompendium2012} in einer umfangreichen Darstellung auf. Nach einer Auseinandersetzung mit Fillmores sowie anderen in und außerhalb der Linguistik vorliegenden Slot-Filler-Frame-Modellen (insbesondere denen von Minsky und Barsalou) schlägt Busse dann ein eigenes Frame-Modell vor. Dieses stellt nicht nur eine Synthese seiner Vorläufer dar, sondern macht Frames und das in ihnen gespeicherte ‚gesamte verstehensrelevante Wissen` auch für die Diskurslinguistik, die er als ‚epistemologische Linguistik` \autocite{busseDiskurslinguistikAlsEpistemologie2008} begreift, zu höchst relevanten Analyseobjekten.

Bereits Mitte des 20. Jahrhunderts finden sich mit den Arbeiten John Rupert Firths und Zellig Harris' außerdem Grundlagen eines ebenfalls völlig neuartigen Zugangs zur Semantik. Ihre Überlegungen zur Ko- und Kontextabhängigkeit von Bedeutung prägen das heute insbesondere im Bereich der computergestützten Sprachverarbeitung weit verbreitete Paradigma der Distributionellen Semantik. Die Grundannahme distributionell-semantischer Ansätze besteht darin, dass die ein Wort umgebenden Wörter Rückschlüsse auf dessen Bedeutung erlauben \autocite[„You shall know a word by the company it keeps.``][11]{firthSynopsisLinguisticTheory1957} und dass demzufolge Wörter, die eine distributionelle Ähnlichkeit aufweisen, d.\,h. von ähnlichen Wörtern umgeben sind, eine ähnliche Bedeutung tragen. Oder anders gewendet: „difference of meaning correlates with difference of distribution`` \autocite[156]{harrisDistributionalStructure1954}.

Diese noch in analogen Sprachanalysen gemachten Beobachtungen sind natürlich gerade für computergestützte Zugänge von Interesse. Für die Korpuslinguistik sind es insbesondere die Arbeiten John Sinclairs und deren Weiterentwicklung durch Elena Tognini-Bonelli, die auf Firths Überlegungen aufbauen, sie in korpuslinguistische Zugänge einbinden und weiterdenken. Die theoretischen Impulse, die sie geben, sind zahlreich. So etwa dass Wörter keinesfalls zufällig miteinander vorkommen oder als Füllungen für grammatische Slots frei wählbar sind (\emph{open-choice principle}), sondern dass sie in gegenseitigen Abhängigkeiten stehen und Wortgrenzen übergreifende Bedeutungseinheiten (\emph{idiom principle}) bilden \autocite{sinclairCorpusConcordanceCollocation1991}. Es sind sowohl die syntagmatische Analyse der Kollokation von Wörtern als auch die paradigmatische Analyse ihrer Verwendungsko(n)texte, die korpuslinguistisch hervorragend zu realisieren und mithin systematisch zu beschreiben sind. Die Rolle korpuslinguistischer Zugänge besteht jedoch in mehr als der bloßen methodischen Ermöglichung dieser Analysen. Sie geht auch mit einer „eigen{[}en{]} Perspektive`` \autocite[2]{perkuhnKorpuslinguistikUnbekannteWesen2006} auf Sprache einher, wie sie insbesondere im von Tognini-Bonelli \autocite*{tognini-bonelliCorpusLinguisticsWork2001} vorgeschlagenen \emph{corpus-driven} Paradigma eingenommen wird. Der Fokus liegt dabei nicht mehr auf der Abfrage zuvor definierter sprachlicher Phänomene oder Hypothesen, sondern der Befragung des Korpus in Hinblick auf sich zeigende Muster und die datengeleitete Gewinnung von Erkenntnissen auf Basis dieser Muster. So können nicht nur größere und mithin repräsentativere Datengrundlagen untersucht, sondern auch introspektiv kaum zugängliche Erkenntnisse gewonnen werden, etwa zu semantischen Prägungen einzelner Wortformen eines Flexionsparadigmas \autocites{sinclairCorpusConcordanceCollocation1991}[92-99]{tognini-bonelliCorpusLinguisticsWork2001} oder zur semantischen Prosodie von Wörtern \autocite[111-116]{tognini-bonelliCorpusLinguisticsWork2001}. Für die korpuslinguistische Diskursanalyse setzt dann Bubenhofer \autocite*{bubenhoferSprachgebrauchsmusterKorpuslinguistikAls2009} auf Basis dieser Theorien -- und anknüpfend an Helmut Feilkes und Angelika Linkes Arbeiten zur von Routinen geprägten sprachlichen Oberfläche bzw. Performanz \autocite{feilkeTextroutineTextsemantikUnd2003, feilkeOberflaecheUndPerformanz2009} -- sprachliche Musterhaftigkeit in Form von ‚Sprachgebrauchsmustern` als bedeutsames Element des Sprachgebrauchs und als Definitionskriterium für Diskurse an. Thematisch ausgerichtete Diskursdefinitionen werden damit abgelehnt und durch \emph{corpus-driven} ermittelte Sprachgebrauchsmuster als „minimale Kristallisationskerne von Diskursen`` \autocite[40]{linkeBegriffsgeschichteDiskursgeschichteSprachgebrauchsgeschichte2003} ersetzt.

Die vorliegende Arbeit unternimmt den Versuch, sich mithilfe datengeleiteter Verfahren das umfangreiche in der Distribution von Wörtern vorliegende Bedeutungs- und Wissenspotenzial zu Nutze zu machen, um semantische Frames im Sinne Busses \autocite*{busseFrameSemantikKompendium2012} für den Extremismusdiskurs herauszuarbeiten. Es wird eine Methode zur Induktion semantischer Frames vorgeschlagen, die mittels der Analyse der Verteilung von Wörtern auf der syntagmatischen und paradigmatischen Achse Frame-semantisch lesbare Wissensbestände an der sprachlichen Oberfläche emergent werden lässt und diese visualisierbar sowie explorierbar macht. Einander entgegengesetzte Theorien bzw. Paradigmen werden dabei aufeinander bezogen und fruchtbar verbunden, nämlich auf der einen Seite die typischerweise qualitativ vorgehende Rekonstruktion von Frames thematisch definierter Diskurse im Sinne einer „Tiefen-Semantik`` \autocite[148]{busseDiskursSpracheGesellschaftliches2013} und auf der anderen Seite die dezidiert nicht nach in der Tiefe, sondern auf der Oberfläche verorteten Wissensbeständen suchende quantitativ-korpuslinguistische Analyse formal definierter Diskurse \autocite{bubenhoferSprachgebrauchsmusterKorpuslinguistikAls2009, scharlothWuchernRhizomeLinguistische2013}. Meines Erachtens profitieren beide Paradigmen von der Inbezugsetzung: Identifizierte sprachoberflächliche Muster können in ein kohärentes und etabliertes semantisches Modell eingebettet werden, während durch den korpuslinguistischen Zugang Probleme der Frame-semantischen empirischen Modellierungsarbeit gelöst werden können. Dies sind insbesondere die kontext- und sprachgebrauchsbasierte Ermittlung einer Vielzahl von Frame-Elementen einer bestimmten Granularität und ihrer wechselseitigen Relationen und lexikalischen Realisierungsmöglichkeiten.

Erprobungsgegenstand für den vorgeschlagenen Ansatz ist der mediale Extremismusdiskurs der Jahre 1999--2021 für dessen Untersuchung ein großes Korpus aus dem Online-Artikelbestand der Zeitungen \textsc{Taz}, \textsc{Spiegel} und \textsc{Welt} zugrunde gelegt wird. Es sollen für unterschiedliche Extremismusformen -- insbesondere Rechts- und Linksextremismus sowie Islamismus -- datengeleitet und kontrastiv Frames und mithin unterschiedliche Formen des Framings rekonstruiert werden. Der Kontrast wird dabei zum einen auf der diachronen Achse angesetzt, indem insgesamt acht Zeitabschnitte untersucht werden; zum anderen sollen die drei Zeitungen in Hinblick auf ihr Framing von einzelnen Extremismusvarianten miteinander verglichen werden.

Die Wahl des Themas Extremismus gründet neben seiner kontinuierlichen Aktualität und gesellschaftlichen Bedeutung auch in der semantischen Unterspezifiziertheit sowie Kontroversität des Extremismusbegriffs. So ist das Konzept in seiner Herkunftsdisziplin, der Politikwissenschaft, hochgradig umstritten und wird mitunter grundsätzlich in Frage gestellt. Von diskursanalytischer Seite wurde außerdem herausgestellt, dass der Extremismusbegriff keinesfalls klar definiert, sondern von unterschiedlichen Akteuren zu unterschiedlicher Zeit unterschiedlich gebraucht wird \autocite{ackermannMetamorphosenExtremismusbegriffesDiskursanalytische2015} und sich einzelne Extremismusvarianten etwa im ihnen zugeschriebenen Gewaltpotenzial deutlich unterscheiden, was ebenfalls zeitlich variabel und von außersprachlichen Ereignissen abhängig ist \autocite{czuloEinflussExtremistischerGewaltereignisse2020}.

Wie also unterschiedliche Medien Extremismus framen, welche Akteure zu welcher Zeit als extremistisch gelten, welche Handlungen sie verüben und wie dies durch Medien bewertet und sprachlich realisiert wird, ist Teil des „gesamten verstehensrelevanten Wissens`` \autocite[805]{busseFrameSemantikKompendium2012}, das es in dieser Arbeit im Rahmen einer epistemologisch ausgerichteten (korpus\mbox{-})linguistischen Diskursanalyse zu erfassen gilt. Genauer lassen sich die verfolgten Forschungsinteressen wie folgt fassen:

\begin{enumerate}
\def\labelenumi{\arabic{enumi}.}
\item
  In Hinblick auf die Verbindung Frame-semantischer und korpuslinguistischer Ansätze soll geprüft werden, ob und inwiefern mithilfe datengeleiteter Korpusverfahren semantische Frames im Sinne Busses \autocite*{busseFrameSemantikKompendium2012} erschlossen werden können.
\item
  Es soll eine Methode entwickelt werden, die geeignet ist, dies auf Basis eines großen Korpus umzusetzen und Frames somit induktiv zu rekonstruieren, visualisieren und kontrastieren. Die ermittelten Wissensstrukturen sollen dabei nicht nur als solche, sondern präzise in ihrem Zusammenwirken mit sprachlichen Einheiten der lexikalischen Ebene beschrieben werden.
\item
  Erprobungsgegenstand für das entwickelte Verfahren ist der Extremismusdiskurs der Jahre 1999--2021 in den Zeitungen \textsc{Taz}, \textsc{Spiegel} und \textsc{Welt}. Es soll das Framing einzelner Extremismusvarianten (insbesondere Rechts- und Linksextremismus sowie Islamismus) erschlossen und kontrastiert werden. Neben Unterschieden zwischen den einzelnen Extremismen sollen auch kontextuelle, d.\,h. diachrone und mediale Faktoren des Framings herausgearbeitet werden.
\end{enumerate}

Zur Bearbeitung dieser Fragen wird das folgende Kapitel zunächst in die theoretischen Grundlagen der Arbeit einführen. Dazu werde ich das Vorhaben in der korpuslinguistisch und epistemologisch ausgerichteten Diskurslinguistik verorten und den zugrundegelegten Diskursbegriff erläutern (\chapref{diskurslinguistik}). Im Anschluss (\chapref{frame-semantik}) werde ich semantische Frames als wesentliches Analyseobjekt dieser auf diskursiv verhandeltes Wissen ausgerichteten Arbeit beleuchten. Dazu wird in aller Kürze die Genese des Fillmoreschen Frame-Konzeptes in der Linguistik nachgezeichnet, um dann die mehrere Frame-Konzepte integrierende Weiterentwicklung des Frame-Begriffs durch Busse \autocite*{busseFrameSemantikKompendium2012} ausführlicher darzustellen. Letztlich gilt es außerdem, mit dem Konzept des ‚Framing` einen in- und außerhalb des wissenschaftlichen Diskurses zirkulierenden Begriff in Hinblick auf seine Eignung für das vorliegende Projekt zu prüfen und zu operationalisieren (\chapref{framing}).

Das darauffolgende \chapref{distributionelle-semantik-korpuspragmatik-und-das-corpus-driven-paradigma} stellt ebenfalls Fragen des diskursiven Wissens in den Fokus, jedoch aus einer gänzlich anderen Perspektive. Anstatt von einem kognitiven Wissensmodell auszugehen, in dem Diskursbeiträge auf Basis diskursiv vermittelten, aber letztlich individuellen Wissens produziert und rezipiert werden, steht hier die bedeutungstragende Musterhaftigkeit der sprachlichen Oberfläche im Vordergrund. Theoretische Zugänge wie die distributionelle Semantik, das \emph{corpus-driven} Paradigma und die Korpuspragmatik zeigen, wie unter Nutzung korpuslinguistischer Methodik reiche Wissensstrukturen aus der musterhaften Verteilung von Wörtern erschlossen und für diskursanalytische Zwecke nutzbar gemacht werden können. Ob man diese sprachoberflächlichen Strukturen als Epiphänomen kognitiven Wissens lesen möchte oder als dezidiert diskursive Wissenseinheiten, wird dabei aber offengelassen.

Nachdem nun die dem linguistisch-methodischen Ansatz dieser Arbeit zugrundeliegenden Theorien aufgearbeitet und ihre potenziellen gegenseitigen Anschlusspunkte freigelegt worden sind, fokussieren die sich anschließenden Theoriekapitel das Extremismuskonzept (\chapref{extremismus-in-den-sozialwissenschaften}) und seine diskursive Aushandlung und mithin das inhaltlich zu bearbeitende Thema der vorliegenden Arbeit. Zunächst werden dazu die sozialwissenschaftlichen Grundlagen des Konzeptes dargelegt sowie die Kritik an ebendiesem geschildert.

Es schließt sich die Darstellung des Forschungsstandes an, in dem zunächst diskursanalytische Arbeiten zum Extremismusdiskurs (\chapref{extremismus-diskurs}) aufgearbeitet werden, die den Begriff als hochgradig veränderlich und -- so Ackermann et al. \autocite*{ackermannMetamorphosenExtremismusbegriffesDiskursanalytische2015} aus einer kritisch-diskursanalytischen Perspektive -- interessengeleitet form- und einsetzbar beschreiben. Neben der wissenschaftlichen Kontextualisierung des vorliegenden Forschungsvorhabens dient das Kapitel auch dazu, grundlegende Charakteristika des Extremismuskonzeptes herauszuarbeiten, die die datengeleitete Analyse informieren können. Zwar soll gerade kein Set an in den Daten zu suchenden Kategorien erarbeitet werden; ein Überblick über zumindest potenziell relevante Wissenselemente ist jedoch notwendig, um überhaupt Strukturen des Extremismusdiskurses in den Korpusdaten erkennen zu können. Daran anschließend werden Ansätze zur Frame-Induktion (\chapref{frame-induktion}) sowie zur korpuslinguistischen Erschließung von Semantik bzw. Diskursen (\chapref{distributionell-semantische-methodik}) dargestellt und in Hinblick auf ihre Potenziale für die zu entwickelnde Methode erörtert. Das Resümee (\chapref{resuemee}) fasst zusammen, warum Frames ein geeignetes Wissensmodell für die Analyse darstellen, inwiefern Frames distributionell-semantisch zu erschließen sind und welche inhaltlichen Charakteristika des Extremismuskonzeptes den Diskurs bislang bestimmt haben und insofern in den Daten zu erwarten sein könnten.

Der dritte Abschnitt der Arbeit stellt das für die Analyse erhobene Korpus vor und arbeitet das methodische Vorgehen aus. Nach der Vorstellung und Begründung der Datengrundlage (\chapref{korpus}) werden zunächst die in der Linguistik variierend verwendeten Begriffe ‚quantitativ` und ‚qualitativ`, \emph{corpus-based} und \emph{corpus-driven} sowie ‚hermeneutisch` und ‚objektiv` definitorisch geschärft und ein Überblick über den Entwicklungsprozess der Methode gegeben (\chapref{prozess-der-methodenentwicklung} und \ref{begriffsklaerungen-quantitativ-und-qualitativ-corpus-driven-und-corpus-based-sowie-die-rolle-von-hermeneutik}). Es werden sodann wesentliche Methoden der distributionellen Semantik, die im vorgeschlagenen Verfahren zum Einsatz kommen, vorgestellt (\chapref{methoden-der-distributionellen-semantik-word-embeddings-und-kollokationen}): zum einen die seit langem in der Korpuslinguistik etablierte Berechnung von Kollokationen, also der Ermittlung syntagmatischer Muster, und zum anderen die vergleichsweise neue und auf der paradigmatischen Achse operierende Berechnung von Word Embeddings. Da Word Embeddings in den Philologien -- ganz im Gegensatz zu den parallelen Entwicklungen in der Computerlinguistik -- bislang kaum eingesetzt wurden, werden sie ausführlicher als Kollokationen eingeführt. Danach erfolgt die Inbezugsetzung und Operationalisierung der mithilfe dieser Methoden freilegbaren sprachlichen (Wissens\mbox{-})Strukturen zum Frame-Begriff Busses (\chapref{operationalisierung-von-frames-fuer-distributionelle-methoden}). Nach diesen theoretischen, die Operationalisierung betreffenden, Überlegungen wird im folgenden \chapref{technisch-methodische-umsetzung-vom-rohen-text-zum-graphen} die technische Seite der entwickelten Methode beleuchtet und der Weg vom ‚rohen` Textkorpus zu visualisierten Frames nachgezeichnet. Ein Kapitel (\ref{wie-der-graph-zu-lesen-ist-und-moegliche-fallstricke}) mit Hinweisen zu Möglichkeiten und Grenzen der hermeneutischen Interpretation der erzeugten Visualisierungen rundet die Ausführungen zu Daten und Methodik ab.

Es schließt sich der empirische Hauptteil der Studie an, der zunächst auf die Makrostruktur der für die einzelnen Teilkorpora ermittelten Frames eingeht (\chapref{makrostruktur-thematische-teil-frames}) und anschließend die für jeden Zeitraum bzw. jede Zeitung modellierten Frames des Rechts- und Linksextremismus sowie des Islamismus beschreibt. Das Analysefazit (\chapref{fazit-der-analyse}) resümiert diese Ergebnisse und stellt darüber hinaus auch varianten- und teilkorpusübergreifende Konstanten des Framings heraus.

In einem letzten Schritt (\chapref{diskussion-limitationen-perspektiven-und-fazit}) werden die Ergebnisse der Analyse mit den eingangs formulierten Zielen abgeglichen und die entwickelte Methode in Hinblick auf ihre Leistungsfähigkeit im Rahmen einer Frame-semantisch ausgerichteten Diskursanalyse diskutiert. Es werden sowohl die inhaltlichen als auch die methodischen Errungenschaften resümiert wie auch verbleibende bzw. neu zu Tage getretene Desiderata für weiterer Forschungsvorhaben benannt.
